% =====================================================================
%  MANUAL DE USUARIO — DEFYVISION METALCONF
%  Sistema de inspeccion visual de chapas perforadas
%  Compilar en Overleaf con: pdfLaTeX
% =====================================================================
\documentclass[11pt,a4paper]{article}

\usepackage[utf8]{inputenc}
\usepackage[T1]{fontenc}
\usepackage[spanish,es-noindentfirst,es-noshorthands]{babel}
\usepackage{lmodern}
\usepackage[margin=2.4cm]{geometry}
\usepackage{xcolor}
\usepackage{graphicx}
\usepackage{booktabs}
\usepackage{array}
\usepackage{enumitem}
\usepackage{tcolorbox}
\usepackage{fancyhdr}
\usepackage{titlesec}
\usepackage{amssymb}
\usepackage[hidelinks]{hyperref}
\tcbuselibrary{skins,breakable}

% ---- Paleta de marca -------------------------------------------------
\definecolor{brand}{HTML}{7C3AED}      % violeta DefyVision
\definecolor{brandDark}{HTML}{4C1D95}
\definecolor{okGreen}{HTML}{166534}
\definecolor{nokRed}{HTML}{991B1B}
\definecolor{warnAmber}{HTML}{B45309}
\definecolor{idleBlue}{HTML}{1D4ED8}
\definecolor{ink}{HTML}{1F2937}
\definecolor{muted}{HTML}{6B7280}
\definecolor{soft}{HTML}{F3F4F6}

% ---- Encabezado / pie ------------------------------------------------
\pagestyle{fancy}
\fancyhf{}
\renewcommand{\headrulewidth}{0.4pt}
\renewcommand{\footrulewidth}{0.4pt}
\lhead{\small\textcolor{brandDark}{\textbf{DEFYVISION METALCONF}}}
\rhead{\small\textcolor{muted}{Manual del operario}}
\lfoot{\small\textcolor{muted}{Inspeccion visual de chapa perforada}}
\rfoot{\small\thepage}

% ---- Titulos de seccion ---------------------------------------------
\titleformat{\section}{\Large\bfseries\color{brandDark}}{\thesection.}{0.6em}{}
\titleformat{\subsection}{\large\bfseries\color{brand}}{\thesubsection}{0.6em}{}

% ---- Cajas reutilizables --------------------------------------------
\newtcolorbox{infobox}{colback=soft,colframe=brand,boxrule=0.8pt,
  arc=2pt,left=8pt,right=8pt,top=6pt,bottom=6pt,breakable}
\newtcolorbox{warnbox}{colback=warnAmber!8,colframe=warnAmber,boxrule=0.8pt,
  arc=2pt,left=8pt,right=8pt,top=6pt,bottom=6pt,breakable,
  title=\textbf{\textcolor{warnAmber}{ATENCION}}}
\newtcolorbox{dangerbox}{colback=nokRed!7,colframe=nokRed,boxrule=0.8pt,
  arc=2pt,left=8pt,right=8pt,top=6pt,bottom=6pt,breakable,
  title=\textbf{\textcolor{nokRed}{IMPORTANTE}}}

\newcommand{\btn}[1]{\colorbox{ink}{\textcolor{white}{\small\,\textbf{#1}\,}}}
\newcommand{\estado}[2]{\colorbox{#1}{\textcolor{white}{\small\,\textbf{#2}\,}}}

% =====================================================================
\begin{document}

% ---------------------- PORTADA --------------------------------------
\begin{titlepage}
\centering
\vspace*{2.5cm}
{\Huge\bfseries\color{brandDark} DEFYVISION METALCONF\par}
\vspace{0.6cm}
{\LARGE\color{brand} Manual del Operario\par}
\vspace{0.3cm}
{\large\color{muted} Sistema de inspeccion visual de chapas perforadas\par}
\vspace{2.2cm}
\begin{tcolorbox}[colback=soft,colframe=brand,width=0.8\textwidth,arc=3pt]
\centering\large
Deteccion automatica de agujeros faltantes\\[4pt]
Comparacion contra patron de referencia\\[4pt]
Clasificacion OK / NOK y parada de maquina\\[4pt]
Integracion con PLC (Coolmay CX3G, Modbus TCP)
\end{tcolorbox}
\vfill
{\large\color{muted} Version del manual: 1.0 \quad$\cdot$\quad Uso interno de planta\par}
\vspace{1cm}
\end{titlepage}

\tableofcontents
\newpage

% =====================================================================
\section{Que hace este sistema}
El programa inspecciona en tiempo real las chapas metalicas perforadas que
pasan por la linea. Con una camara y una luz de fondo (\emph{backlight}),
detecta los agujeros de cada pieza y los compara contra un \textbf{patron de
referencia} previamente calibrado. Si a una pieza le faltan agujeros, el
sistema la marca como defectuosa y, si el problema persiste, \textbf{detiene
la maquina} para que se revise.

\begin{infobox}
\textbf{En una frase:} el sistema mira cada chapa, cuenta y ubica los
agujeros, y avisa (o para la maquina) cuando encuentra un defecto real y
sostenido.
\end{infobox}

El equipo trabaja con \textbf{dos scanners} independientes (Scanner 1 y
Scanner 2), cada uno con su propia camara, su patron y su semaforo.

% =====================================================================
\section{Encendido y arranque del programa}
\begin{enumerate}[leftmargin=1.4em]
  \item Encender la PC, el PLC, las camaras y la iluminacion de fondo.
  \item Abrir el programa desde el acceso directo del escritorio
        (\textbf{DEFYVISION}) o, si se trabaja por consola, con el comando
        de \emph{produccion}.
  \item Esperar a que aparezca la pantalla de operario con los dos scanners
        y que el indicador de \textbf{PLC} figure como \textcolor{okGreen}{Conectado}.
\end{enumerate}

\begin{infobox}
\textbf{Tres formas de abrir el sistema:}
\begin{itemize}[leftmargin=1.2em,nosep]
  \item \textbf{Produccion} (\texttt{run}): PLC + camaras + pantalla de
        operario en vivo. Es el modo normal de trabajo.
  \item \textbf{Servicio} (\texttt{service}): pantalla tecnica con
        \textbf{clave de acceso}. Solo personal autorizado (calibracion,
        diagnostico, ROI, tolerancias).
  \item \textbf{Analisis de carpeta} (\texttt{operator-ui}): revisar
        imagenes ya guardadas, sin camara ni PLC.
\end{itemize}
\end{infobox}

\begin{dangerbox}
Lo mas importante: \textbf{la aplicacion siempre debe poder abrirse}. Si al
arrancar aparece un error de consola y no abre la ventana, no insista:
avise a soporte tecnico. Nunca modifique archivos de configuracion sin
autorizacion.
\end{dangerbox}

% =====================================================================
\section{La pantalla del operario}
Cada scanner ocupa un panel con la siguiente informacion:

\begin{center}
\renewcommand{\arraystretch}{1.35}
\begin{tabular}{>{\raggedright\arraybackslash}p{4.6cm} p{10.2cm}}
\toprule
\textbf{Elemento} & \textbf{Que significa}\\
\midrule
\textbf{Titulo} & \texttt{SCANNER 1 $\cdot$ Microperforado} (numero de scanner y modelo activo).\\
\textbf{Estado} & Cartel de color con el estado actual (ver Seccion~\ref{sec:estados}).\\
\textbf{Modo} & \texttt{MANUAL} o \texttt{AUTO} (posicion de la maneta).\\
\textbf{Racha} & Frames NOK seguidos / umbral para parar (ej. \texttt{2 / 4}).\\
\textbf{NOK} & Cantidad de piezas defectuosas de la sesion.\\
\textbf{OK} & Cantidad de piezas correctas de la sesion.\\
\textbf{Tiempo} & Tiempo en marcha desde que se apreto INICIAR.\\
\textbf{Camara} & Vista en vivo con borde verde (en marcha) o rojo (detenido).\\
\bottomrule
\end{tabular}
\end{center}

% =====================================================================
\section{Estados del sistema}\label{sec:estados}
El cartel de estado cambia de color y texto segun lo que esta pasando:

\begin{center}
\renewcommand{\arraystretch}{1.5}
\begin{tabular}{c >{\raggedright\arraybackslash}p{11.5cm}}
\toprule
\textbf{Cartel} & \textbf{Significado y que hacer}\\
\midrule
\estado{idleBlue}{EN ESPERA} & El scanner esta listo pero no inspeccionando.
  Apriete \btn{INICIAR} para comenzar.\\
\estado{okGreen}{INSPECCIONANDO} & Funcionando normal, analizando cada pieza.
  No requiere accion.\\
\estado{warnAmber}{FALLA DETECTADA} & Se detuvo por una racha de piezas
  defectuosas. Revise la pieza/punzon y haga \btn{RESET FALLA}.\\
\estado{nokRed}{DETENIDO} & La maquina fue detenida (por falla o manualmente).
  Revise y reinicie el ciclo.\\
\estado{ink}{ERROR} & Problema tecnico (ej. camara sin senal). Ver
  Seccion~\ref{sec:errores}.\\
\bottomrule
\end{tabular}
\end{center}

\begin{infobox}
\textbf{Diferencia clave:} \estado{warnAmber}{FALLA DETECTADA} y
\estado{nokRed}{DETENIDO} ambos frenan la produccion, pero la
\textbf{PARADA DE MAQUINA} (Seccion~\ref{sec:parada}) ademas muestra en
pantalla completa la \emph{foto de la pieza} que causo el problema.
\end{infobox}

% =====================================================================
\section{Modos: MANUAL y AUTO}
El modo lo define la \textbf{maneta} conectada al PLC:
\begin{itemize}[leftmargin=1.4em]
  \item \textbf{AUTO:} el sistema inspecciona y, ante un defecto sostenido,
        acciona la salida hacia el PLC (corta la electrovalvula / detiene).
        Es el modo de produccion.
  \item \textbf{MANUAL:} el sistema inspecciona y muestra resultados, pero
        no acciona la parada automatica. Util para pruebas y ajuste.
\end{itemize}

\begin{warnbox}
Si el cartel dice \texttt{MANETA: SIN SENAL}, el PLC no esta leyendo la
posicion de la maneta. Avise a mantenimiento antes de producir.
\end{warnbox}

% =====================================================================
\section{Botones y flujo de trabajo}
\begin{center}
\renewcommand{\arraystretch}{1.4}
\begin{tabular}{c >{\raggedright\arraybackslash}p{11.8cm}}
\toprule
\textbf{Boton} & \textbf{Accion}\\
\midrule
\btn{$\blacktriangleright$ INICIAR} & Arranca la inspeccion. El estado pasa a
  \estado{okGreen}{INSPECCIONANDO}.\\
\btn{$\blacksquare$ DETENER} & Detiene la inspeccion de forma controlada.\\
\btn{$\circlearrowleft$ RESET FALLA} & Aparece solo cuando el scanner se detuvo
  por una falla. Limpia la falla y deja el scanner \estado{idleBlue}{EN ESPERA}.\\
\bottomrule
\end{tabular}
\end{center}

\textbf{Ciclo normal de trabajo:}
\begin{enumerate}[leftmargin=1.5em]
  \item Verificar \textcolor{okGreen}{PLC conectado} y modo \texttt{AUTO}.
  \item Apretar \btn{INICIAR} en cada scanner.
  \item Trabajar. El sistema cuenta OK/NOK solo.
  \item Ante una parada por falla: revisar la pieza $\to$ \btn{RESET FALLA}
        $\to$ \btn{INICIAR} de nuevo.
  \item Al terminar el turno: \btn{DETENER}.
\end{enumerate}

% =====================================================================
\section{El semaforo (luces del PLC)}
Cada scanner comanda una torre de luces:

\begin{center}
\renewcommand{\arraystretch}{1.4}
\begin{tabular}{c >{\raggedright\arraybackslash}p{11.5cm}}
\toprule
\textbf{Luz} & \textbf{Significado}\\
\midrule
\textcolor{idleBlue}{$\blacksquare$} \textbf{Azul} & En espera / listo para arrancar.\\
\textcolor{okGreen}{$\blacksquare$} \textbf{Verde} & Funcionando correctamente.\\
\textcolor{warnAmber}{$\blacksquare$} \textbf{Amarilla} & Racha de NOK en curso (advertencia; aun no paro).\\
\textcolor{nokRed}{$\blacksquare$} \textbf{Roja} & Falla / parada: requiere intervencion.\\
\bottomrule
\end{tabular}
\end{center}

\begin{infobox}
La \textbf{luz de fondo (backlight)} de la inspeccion permanece \emph{siempre
encendida} mientras el sistema esta activo; no es parte del semaforo.
\end{infobox}

% =====================================================================
\section{Parada de maquina y fallas}\label{sec:parada}
El sistema \textbf{no para por una sola pieza dudosa}. Para evitar falsas
alarmas, exige que el defecto se repita durante \textbf{varios frames
seguidos} antes de detener:

\begin{center}
\renewcommand{\arraystretch}{1.4}
\begin{tabular}{>{\raggedright\arraybackslash}p{5.2cm} >{\raggedright\arraybackslash}p{9.6cm}}
\toprule
\textbf{Tipo de parada} & \textbf{Como se dispara}\\
\midrule
\textbf{Faltantes persistentes} & Una misma columna/zona con agujeros
  ausentes durante varios frames consecutivos (probable punzon roto o tapado).\\
\textbf{Patron desalineado} & La chapa entra corrida o inclinada de forma
  sostenida (problema mecanico), no un agujero suelto.\\
\bottomrule
\end{tabular}
\end{center}

\begin{infobox}
\textbf{Regla de los frames (Scanner 1):} ninguna parada se dispara antes de
\textbf{4 frames} consecutivos con el mismo problema. Scanner 2 usa su propio
ajuste. Con la chapa detenida los frames son casi identicos, por eso la
parada puede sentirse ``instantanea'', pero siempre requiere varios frames.
\end{infobox}

\textbf{Cuando ocurre una PARADA DE MAQUINA:}
\begin{enumerate}[leftmargin=1.5em]
  \item Aparece una \textbf{alerta a pantalla completa} con la foto de la
        pieza que causo la parada y el motivo.
  \item La luz roja queda encendida y la electrovalvula se corta.
  \item El operario debe: \textbf{retirar/revisar la pieza} $\to$ apretar
        \btn{CONFIRMAR Y CONTINUAR} (o \btn{RESET FALLA}) $\to$ \btn{INICIAR}.
\end{enumerate}

\begin{warnbox}
Si la maquina para muy seguido \emph{sin un defecto real visible}, casi
siempre es un problema de \textbf{calibracion} (iluminacion, suciedad en el
lente, ROI corrido o chapa mal posicionada), no un problema del programa.
Avise a personal de servicio para recalibrar.
\end{warnbox}

% =====================================================================
\section{Tabla de errores frecuentes}\label{sec:errores}
\begin{center}
\renewcommand{\arraystretch}{1.45}
\begin{tabular}{>{\raggedright\arraybackslash}p{4.8cm} >{\raggedright\arraybackslash}p{10cm}}
\toprule
\textbf{Sintoma} & \textbf{Causa probable / solucion}\\
\midrule
\estado{ink}{ERROR} \textbf{o} ``CAMARA'' con segundos & La camara dejo de
  enviar imagen. Revisar cable de red / alimentacion de la camara. Reiniciar
  si persiste.\\
\texttt{PLC: Desconectado} & Sin comunicacion con el PLC. Revisar red y que
  el PLC este encendido.\\
\texttt{MANETA: SIN SENAL} & El PLC no lee el selector MANUAL/AUTO. Avisar a
  mantenimiento.\\
Muchas piezas NOK siendo buenas & Iluminacion cambiada, lente sucio o ROI
  corrido. Limpiar y, si sigue, recalibrar (modo Servicio).\\
Parada de maquina muy frecuente & Punzon realmente roto/tapado \emph{o}
  desalineacion mecanica de la chapa. Revisar primero la pieza real.\\
La ventana no abre / consola con texto rojo & Problema tecnico de arranque.
  No tocar archivos; llamar a soporte.\\
\bottomrule
\end{tabular}
\end{center}

% =====================================================================
\section{Modo Servicio (personal autorizado)}
Se abre con el comando \texttt{service} e ingresando la \textbf{clave} en la
pantalla ``Activar sistema''. Da acceso a herramientas tecnicas:

\begin{itemize}[leftmargin=1.4em]
  \item \textbf{GRABACION / ANALISIS / CALIBRACION / CONEXION:} vistas de
        camara para grabar evidencia, analizar y ajustar la toma.
  \item \textbf{PLC I/O:} estado de entradas/salidas del PLC en vivo.
  \item \textbf{Diagnostico:} informacion tecnica del sistema.
  \item \textbf{Tolerancias:} parametros de deteccion (umbral, area, etc.).
  \item \textbf{Area de analisis:} incluye el \textbf{editor de ROI}.
\end{itemize}

\begin{dangerbox}
Las \textbf{Tolerancias} y el \textbf{ROI} definen si una pieza es OK o NOK.
Modificarlos mal puede llenar de falsas alarmas o, peor, dejar pasar piezas
defectuosas. Solo debe tocarlos personal capacitado, y siempre validando con
piezas conocidas.
\end{dangerbox}

% =====================================================================
\section{Editor de ROI (region de inspeccion)}
El \textbf{ROI} es el rectangulo dentro de la imagen donde el sistema busca
los agujeros (la franja de la chapa). Se ajusta en \emph{Area de analisis} o
en el \emph{modo Servicio}.

\begin{infobox}
\textbf{Ancho fijo:} el ROI tiene un \textbf{ancho constante}. Solo se puede
\textbf{mover a izquierda o derecha} con los botones \btn{MOVER ROI $\blacktriangleleft$}
\ / \btn{$\blacktriangleright$}. Ya \textbf{no} se puede hacer mas ancho ni mas
angosto (eso causaba errores). El campo \texttt{PASO px} define cuantos pixeles
se mueve por click.
\end{infobox}

\textbf{Para reubicar el ROI:}
\begin{enumerate}[leftmargin=1.5em]
  \item Abrir una imagen de referencia o activar \btn{Camara en vivo}.
  \item Con \btn{MOVER ROI} centrar el rectangulo sobre la chapa.
  \item Apretar \btn{GUARDAR ROI}. El sistema guarda la posicion \textbf{y
        reconstruye el patron} automaticamente.
\end{enumerate}

\begin{warnbox}
La posicion del ROI \textbf{queda guardada y no se mueve sola} al reiniciar el
programa. Si notas que ``se corrio'', es porque alguien lo movio manualmente;
volve a centrarlo y \btn{GUARDAR ROI}.
\end{warnbox}

% =====================================================================
\section{Buenas practicas y mantenimiento}
\begin{itemize}[leftmargin=1.4em]
  \item Mantener el \textbf{lente limpio} y la \textbf{luz de fondo} sin
        obstrucciones ni cambios de intensidad.
  \item Verificar que la chapa entre \textbf{derecha y centrada} en el ROI.
  \item Al empezar el turno, controlar contadores OK/NOK en cero y estado
        \estado{idleBlue}{EN ESPERA}.
  \item Ante cambios de pieza, lente, camara o iluminacion, \textbf{recalibrar}
        (patron + tolerancias) con personal de servicio.
  \item No apagar la PC de golpe; cerrar el programa correctamente.
\end{itemize}

% =====================================================================
\section{Glosario}
\begin{description}[leftmargin=2.2cm,style=nextline]
  \item[ROI] Region de interes: la franja de la imagen donde se inspecciona.
  \item[Patron] Mapa de referencia de donde deben estar los agujeros de una
        pieza buena.
  \item[OK / NOK] Pieza correcta / pieza defectuosa (le faltan agujeros).
  \item[Frame] Cada cuadro de imagen que analiza la camara.
  \item[Racha NOK] Cantidad de frames defectuosos seguidos.
  \item[Parada de maquina] Detencion automatica por defecto sostenido.
  \item[Backlight] Luz de fondo que ilumina la chapa por detras.
  \item[PLC] Controlador que maneja luces, electrovalvulas y sensores.
  \item[Modbus TCP] Protocolo de comunicacion con el PLC por red.
\end{description}

% =====================================================================
\section*{Soporte tecnico}
\addcontentsline{toc}{section}{Soporte tecnico}
\begin{infobox}
Ante cualquier duda, error que no figure en este manual, o comportamiento
inesperado del sistema, \textbf{no modifique la configuracion}: registre lo
que paso (foto de la pantalla si es posible) y contacte a soporte tecnico.
\end{infobox}

\end{document}
